% GENERATED FILE. Edit canonical lessons, then run npm run generate:books.
% canonical-source-hash: fnv1a64:3e2659b838da5d83
% canonical-lessons: TE-C500-anandincu, TE-C500-murisipovu, TE-C500-pongipovu, TE-C500-pulakincu, TE-C500-ullasam

\chapter{Joy and Delight}
\label{ch:te-joy-and-delight}
\hlchaptermodality{500}

\begin{chapteropening}
\textbf{By the end of this chapter:} \emph{I can say that I enjoyed something or was delighted.}

Complete the last lesson of chapter 500: I can say that I enjoyed something or was delighted.
\end{chapteropening}

\section[\texorpdfstring{\te{ఆనందించు} (ānandiñcu)}{ఆనందించు (ānandiñcu)}]{\te{ఆనందించు} (ānandiñcu) — to enjoy; to rejoice}

\label{lesson:TE-C500-anandincu}



\begin{quote}
Before the new one: say the Telugu for a grudge, then the Telugu for contempt.
\end{quote}

\subsection*{You'll want to know: \te{ఆనందించు}}
\textbf{\te{ఆనందించు}} — \emph{ānandiñcu} — \textquotedblleft{}to enjoy; to rejoice\textquotedblright{}.

\begin{quote}
\te{సెలవుల్లో} \te{పిల్లలు} \te{సముద్రం} \te{దగ్గర} \te{ఆనందించారు}.
\end{quote}

\begin{quote}
\emph{selavullō pillalu samudraṁ daggara ānandiñcāru.}
\end{quote}

\begin{quote}
In the holidays the children enjoyed themselves by the sea.
\end{quote}

Past (I …): \te{ఆనందించాను}, \emph{ānandiñcānu}.

The noun is \te{ఆనందం}, joy; \te{ఆనందంగా}, happily.

\begin{cousinweb}
\textbf{\te{ఆనందించు}} (\emph{ānandiñcu}) is built on Sanskrit \emph{ānanda}, \textquotedblleft{}happiness\textquotedblright{} (CDIAL 1172), from the root \emph{nand}, with the Telugu verb ending -\te{ఇంచు}.
\end{cousinweb}

\begin{grammarlens}[title={What you've built}]
One more word for describing what you lived and felt.
\end{grammarlens}

\subsection*{Your turn}
\begin{itemize}
\raggedright
  \item \emph{Say it:} \emph{ānandiñcu}
  \item \emph{Say it:} \emph{ānandiñcu}, once more
  \item \emph{Say it:} \emph{selavullō pillalu samudraṁ daggara ānandiñcāru.} — In the holidays the children enjoyed themselves by the sea.
  \item \emph{Recall:} say the Telugu for a grudge, then the Telugu for contempt, then say \emph{ānandiñcu} again
\end{itemize}

\begin{tcolorbox}[breakable,colback=teal!4,colframe=teal!35!black,title={Before you move on}]
What does \te{ఆనందించు} mean? (\textquotedblleft{}To enjoy; to rejoice\textquotedblright{}.) Say it once more.
\end{tcolorbox}
\section[\texorpdfstring{\te{మురిసిపోవు} (murisipōvu)}{మురిసిపోవు (murisipōvu)}]{\te{మురిసిపోవు} (murisipōvu) — to be delighted}

\label{lesson:TE-C500-murisipovu}



\begin{quote}
Before the new one: say the Telugu for contempt, then the Telugu for to enjoy.
\end{quote}

\subsection*{You'll want to know: \te{మురిసిపోవు}}
\textbf{\te{మురిసిపోవు}} — \emph{murisipōvu} — \textquotedblleft{}to be delighted\textquotedblright{}.

\begin{quote}
\te{మనవడిని} \te{చూసి} \te{తాతయ్య} \te{మురిసిపోయారు}.
\end{quote}

\begin{quote}
\emph{manavaḍini cūsi tātayya murisipōyāru.}
\end{quote}

\begin{quote}
Seeing his grandson, grandpa was delighted.
\end{quote}

Past (I …): \te{మురిసిపోయాను}, \emph{murisipōyānu}.

A warm word for the glow of a parent or a grandparent.

\begin{cousinweb}
\textbf{\te{మురిసిపోవు}} (\emph{murisipōvu}) is native Dravidian: DEDR 5011 lists Telugu \emph{muriyu}, \textquotedblleft{}to exult, to be in great glee\textquotedblright{}, and the verb ends in \emph{pōvu}, \textquotedblleft{}to go\textquotedblright{} (DEDR 4572).
\end{cousinweb}

\begin{grammarlens}[title={What you've built}]
One more word for describing what you lived and felt.
\end{grammarlens}

\subsection*{Your turn}
\begin{itemize}
\raggedright
  \item \emph{Say it:} \emph{murisipōvu}
  \item \emph{Say it:} \emph{murisipōvu}, once more
  \item \emph{Say it:} \emph{manavaḍini cūsi tātayya murisipōyāru.} — Seeing his grandson, grandpa was delighted.
  \item \emph{Recall:} say the Telugu for contempt, then the Telugu for to enjoy, then say \emph{murisipōvu} again
\end{itemize}

\begin{tcolorbox}[breakable,colback=teal!4,colframe=teal!35!black,title={Before you move on}]
What does \te{మురిసిపోవు} mean? (\textquotedblleft{}To be delighted\textquotedblright{}.) Say it once more.
\end{tcolorbox}
\section[\texorpdfstring{\te{పొంగిపోవు} (poṅgipōvu)}{పొంగిపోవు (poṅgipōvu)}]{\te{పొంగిపోవు} (poṅgipōvu) — to be overjoyed; to swell with pride}

\label{lesson:TE-C500-pongipovu}



\begin{quote}
Before the new one: say the Telugu for to enjoy, then the Telugu for to be delighted.
\end{quote}

\subsection*{You'll want to know: \te{పొంగిపోవు}}
\textbf{\te{పొంగిపోవు}} — \emph{poṅgipōvu} — \textquotedblleft{}to be overjoyed; to swell with pride\textquotedblright{}.

\begin{quote}
\te{బహుమతి} \te{గెలిచినప్పుడు} \te{ఆమె} \te{పొంగిపోయింది}.
\end{quote}

\begin{quote}
\emph{bahumati gelicinappuḍu āme poṅgipōyindi.}
\end{quote}

\begin{quote}
When she won the prize, she was overjoyed.
\end{quote}

Past (I …): \te{పొంగిపోయాను}, \emph{poṅgipōyānu}.

\te{పొంగు} is to boil over, like milk: joy that overflows.

\begin{cousinweb}
\textbf{\te{పొంగిపోవు}} (\emph{poṅgipōvu}) is native Dravidian: DEDR 4469 lists Telugu \emph{poṅgu}, \textquotedblleft{}to bubble up, to boil; to rejoice, to be elated\textquotedblright{}, with Tamil \emph{poṅku}; \emph{pōvu}, \textquotedblleft{}to go\textquotedblright{} (DEDR 4572) ends it.
\end{cousinweb}

\begin{grammarlens}[title={What you've built}]
One more word for describing what you lived and felt.
\end{grammarlens}

\subsection*{Your turn}
\begin{itemize}
\raggedright
  \item \emph{Say it:} \emph{poṅgipōvu}
  \item \emph{Say it:} \emph{poṅgipōvu}, once more
  \item \emph{Say it:} \emph{bahumati gelicinappuḍu āme poṅgipōyindi.} — When she won the prize, she was overjoyed.
  \item \emph{Recall:} say the Telugu for to enjoy, then the Telugu for to be delighted, then say \emph{poṅgipōvu} again
\end{itemize}

\begin{tcolorbox}[breakable,colback=teal!4,colframe=teal!35!black,title={Before you move on}]
What does \te{పొంగిపోవు} mean? (\textquotedblleft{}To be overjoyed; to swell with pride\textquotedblright{}.) Say it once more.
\end{tcolorbox}
\section[\texorpdfstring{\te{పులకించు} (pulakiñcu)}{పులకించు (pulakiñcu)}]{\te{పులకించు} (pulakiñcu) — to thrill; to be thrilled}

\label{lesson:TE-C500-pulakincu}



\begin{quote}
Before the new one: say the Telugu for to be delighted, then the Telugu for to be overjoyed.
\end{quote}

\subsection*{You'll want to know: \te{పులకించు}}
\textbf{\te{పులకించు}} — \emph{pulakiñcu} — \textquotedblleft{}to thrill; to be thrilled\textquotedblright{}.

\begin{quote}
\te{ఆ} \te{పాట} \te{విని} \te{నా} \te{మనసు} \te{పులకించింది}.
\end{quote}

\begin{quote}
\emph{ā pāṭa vini nā manasu pulakiñcindi.}
\end{quote}

\begin{quote}
Hearing that song, my heart thrilled.
\end{quote}

Past (I …): \te{పులకించాను}, \emph{pulakiñcānu}.

Said of the body or the heart: the hair stands on end with joy.

\begin{cousinweb}
\textbf{\te{పులకించు}} (\emph{pulakiñcu}) is built on Sanskrit \emph{pulaka}, the hair of the body standing on end, a sign of rapture rather than fear (Monier-Williams); the Telugu ending -\te{ఇంచు} makes it a verb.
\end{cousinweb}

\begin{grammarlens}[title={What you've built}]
One more word for describing what you lived and felt.
\end{grammarlens}

\subsection*{Your turn}
\begin{itemize}
\raggedright
  \item \emph{Say it:} \emph{pulakiñcu}
  \item \emph{Say it:} \emph{pulakiñcu}, once more
  \item \emph{Say it:} \emph{ā pāṭa vini nā manasu pulakiñcindi.} — Hearing that song, my heart thrilled.
  \item \emph{Recall:} say the Telugu for to be delighted, then the Telugu for to be overjoyed, then say \emph{pulakiñcu} again
\end{itemize}

\begin{tcolorbox}[breakable,colback=teal!4,colframe=teal!35!black,title={Before you move on}]
What does \te{పులకించు} mean? (\textquotedblleft{}To thrill; to be thrilled\textquotedblright{}.) Say it once more.
\end{tcolorbox}
\section[\texorpdfstring{\te{ఉల్లాసం} (ullāsaṁ)}{ఉల్లాసం (ullāsaṁ)}]{\te{ఉల్లాసం} (ullāsaṁ) — cheerfulness; gaiety}

\label{lesson:TE-C500-ullasam}



\begin{quote}
Before the new one: say the Telugu for to be overjoyed, then the Telugu for to thrill.
\end{quote}

\subsection*{You'll want to know: \te{ఉల్లాసం}}
\textbf{\te{ఉల్లాసం}} — \emph{ullāsaṁ} — \textquotedblleft{}cheerfulness; gaiety\textquotedblright{}.

\begin{quote}
\te{పండుగ} \te{రోజు} \te{ఇల్లంతా} \te{ఉల్లాసంగా} \te{ఉంది}.
\end{quote}

\begin{quote}
\emph{paṇḍuga rōju illantā ullāsaṅgā undi.}
\end{quote}

\begin{quote}
On the festival day the whole house is cheerful.
\end{quote}

\te{ఉల్లాసంగా}, cheerfully, in high spirits.

\begin{cousinweb}
\textbf{\te{ఉల్లాసం}} (\emph{ullāsaṁ}) is borrowed from Sanskrit \emph{ullāsa}, both \textquotedblleft{}light\textquotedblright{} and \textquotedblleft{}joy, merriness\textquotedblright{} (Monier-Williams).
\end{cousinweb}

\begin{grammarlens}[title={What you've built}]
One more word for describing what you lived and felt.
\end{grammarlens}

\subsection*{Your turn}
\begin{itemize}
\raggedright
  \item \emph{Say it:} \emph{ullāsaṁ}
  \item \emph{Say it:} \emph{ullāsaṁ}, once more
  \item \emph{Say it:} \emph{paṇḍuga rōju illantā ullāsaṅgā undi.} — On the festival day the whole house is cheerful.
  \item \emph{Recall:} say the Telugu for to be overjoyed, then the Telugu for to thrill, then say \emph{ullāsaṁ} again
\end{itemize}

\begin{tcolorbox}[breakable,colback=teal!4,colframe=teal!35!black,title={Before you move on}]
What does \te{ఉల్లాసం} mean? (\textquotedblleft{}Cheerfulness; gaiety\textquotedblright{}.) Say it once more.
\end{tcolorbox}
